\DomainHeader{Customizing USD}{6\%}

\begin{Objectives}
  \item Build a USD plug-in against a specific USD version.
  \item Build USD from source with custom dependencies.
  \item Create custom model kinds when appropriate.
  \item Create custom schemas for proprietary data models.
  \item Integrate a custom resolver for dynamic asset paths.
  \item Use schemas for nonstandard attributes/structures during import/export.
  \item Write a SceneIndex plug-in that generates Hydra prims.
  \item Write an AssetResolver that generates in-memory renderable primitives.
\end{Objectives}

\begin{Resources}
  \item Learn OpenUSD: Understanding Model Kinds
  \item OpenUSD tutorial: Generating New Schema Classes (usdGenSchema)
  \item Whitepaper: What Are OpenUSD Schemas?
  \item OpenUSD API: PlugRegistry; Kind --- Extensible Categorization;
        Ar (Asset Resolution); file format plug-ins
\end{Resources}

\CheatSheetTitle
\begin{itemize}[leftmargin=*]
  \item \textbf{Everything pluggable goes through one door}:
        \texttt{plugInfo.json} files on \texttt{PXR\_PLUGINPATH\_NAME},
        loaded by \texttt{PlugRegistry} --- schemas, kinds, file formats,
        resolvers, Hydra plug-ins.\index{plugInfo.json}\index{PlugRegistry}
  \item \textbf{Typed (IsA) schema}: defines a new prim type
        (\texttt{def SimJoint}). \textbf{API schema}: adds properties and
        behavior to \emph{existing} types (applied; recorded in
        \texttt{apiSchemas} metadata).\index{IsA schema}\index{API schema}
  \item \textbf{usdGenSchema} turns a \texttt{schema.usda} description into
        registered schemas; with \texttt{skipCodeGeneration = true} they are
        \emph{codeless} --- deployable without compiling.\index{usdGenSchema}\index{codeless schema}
  \item \textbf{The formality ladder}: ad-hoc custom attributes $\to$
        namespaced conventions $\to$ API schema $\to$ typed schema. Climb
        only when querying/validation/tooling demands it.
  \item \textbf{Custom kinds} extend the model hierarchy vocabulary via
        plugInfo (\texttt{"vehicle" baseKind "component"}) --- every site
        that opens your assets needs the same registration.\index{kind}
  \item \textbf{Ar resolver}: logical identifier $\to$ resolved path;
        contexts pin versions/sites. \textbf{File format plug-in}: teaches
        USD to read a foreign format \emph{as} scene description.
        \textbf{SceneIndex}: procedurally feeds Hydra prims at render time.
  \item \textbf{Version pinning}: plug-ins link against one exact USD
        build --- ABI compatibility is per-version, so match headers,
        compiler, and flags.
\end{itemize}

\section{Core Concepts}

\subsection{The plugin map}\label{sec:cu-map}\objref{3.1}

\begin{figure}[ht]
  \centering
  \begin{tikzpicture}[node distance=5mm and 8mm]
    \node[layerbox, minimum width=3.4cm] (reg) {PlugRegistry};
    \node[notetext, above=1mm of reg] {discovers \texttt{plugInfo.json} on \texttt{PXR\_PLUGINPATH\_NAME}};
    \node[primbox, below left=9mm and 14mm of reg] (schemas) {schemas};
    \node[primbox, below left=9mm and -8mm of reg] (kinds) {kinds};
    \node[primbox, below right=9mm and -8mm of reg] (ar) {Ar resolvers};
    \node[primbox, below right=9mm and 14mm of reg] (ff) {file formats};
    \node[primbox, below=20mm of reg] (hydra) {Hydra / SceneIndex};
    \draw[arc, draw=black!55] (reg.south) -- (schemas.north);
    \draw[arc, draw=black!55] (reg.south) -- (kinds.north);
    \draw[arc, draw=black!55] (reg.south) -- (ar.north);
    \draw[arc, draw=black!55] (reg.south) -- (ff.north);
    \draw[arc, draw=black!55] (reg.south) -- (hydra.north);
  \end{tikzpicture}
  \caption{One discovery mechanism, five common extension points. If a
  plug-in ``doesn't load'', debug the door first:
  \texttt{TF\_DEBUG=PLUG\_REGISTRATION}.}
  \label{fig:plugin-map}
\end{figure}

\subsection{Custom schemas: describe once, generate the rest}\label{sec:cu-schemas}\objref{3.4, 3.6}
\index{usdGenSchema}A proprietary data model becomes a first-class USD
citizen through a \texttt{schema.usda} description that
\texttt{usdGenSchema} consumes:

\begin{lstlisting}
#usda 1.0
(
    subLayers = [@usd/schema.usda@]
)
over "GLOBAL" (
    customData = {
        string libraryName = "simRobot"
        string libraryPath = "."
        bool skipCodeGeneration = true
    }
)
{
}
class SimJoint "SimJoint" (
    inherits = </Typed>
)
{
    float torqueLimit = 100 (
        doc = "Maximum torque in Nm."
    )
}
class "WearAPI" (
    inherits = </APISchemaBase>
    customData = {
        token apiSchemaType = "singleApply"
    }
)
{
    float wearAmount = 0 (
        doc = "0 = pristine, 1 = destroyed."
    )
}
\end{lstlisting}

\noindent One file declares both flavors: \texttt{SimJoint} is a
\emph{typed} (IsA) schema --- prims can be \texttt{def SimJoint} --- and
\texttt{WearAPI} is an \emph{applied API} schema that any existing prim type
can opt into. With \texttt{skipCodeGeneration = true} the result is a
\emph{codeless} schema: registered purely as data (plugInfo +
generated schema layer), no C++ to compile or deploy --- the modern default
for pipeline schemas.\index{codeless schema}

\begin{ExamTip}
Typed vs API in one breath: \emph{is it a thing, or a capability of existing
things?} New prim type with its own identity $\Rightarrow$ typed/IsA.
Extra properties for prims that already have a type (physics on meshes,
wear on props) $\Rightarrow$ applied API schema. Reach for schemas at all
only when ad-hoc namespaced attributes stop being enough --- fallbacks,
docs, validation, generated accessors.
\end{ExamTip}

\subsection{Applied API schemas in action}\label{sec:cu-api}\objref{3.4}
Application is recorded on the prim (in \texttt{apiSchemas} metadata) and
queryable --- here with a built-in multiple-apply schema, which behaves
exactly like the ones usdGenSchema generates for you:

\begin{lstlisting}[language=Python]
from pxr import Usd, UsdGeom

stage = Usd.Stage.CreateInMemory()
prim = UsdGeom.Xform.Define(stage, "/Robot").GetPrim()

Usd.CollectionAPI.Apply(prim, "render")
print(prim.HasAPI(Usd.CollectionAPI))   # True
print(prim.GetAppliedSchemas())         # ['CollectionAPI:render']
\end{lstlisting}

\noindent\texttt{Apply()} records the intent, \texttt{HasAPI()} answers
cheaply, and multiple-apply schemas take an instance name (here
\texttt{render}) so one capability can exist several times on a prim.

\subsection{Building USD itself}\label{sec:cu-build}\objref{3.2, 3.1}
\index{build\_usd.py}Sometimes the pipeline needs its own USD: a DCC ships
one build, the farm needs another, or a plug-in requires components the
prebuilt binaries exclude. The supported path is the repository's
\texttt{build\_usd.py}, which fetches dependencies and drives CMake; the
decisions that matter are the flags --- \texttt{--no-imaging} (core only, no
Hydra) for headless services, \texttt{--no-python} for embedded use,
optional components like MaterialX or OpenVDB, and monolithic vs shared
libraries. ``Custom dependencies'' in practice means pinning versions ---
TBB, MaterialX, OpenSubdiv, and above all the \emph{compiler and flags},
because every plug-in must be built with the same toolchain as the USD that
loads it.

\begin{ExamTip}
Exam framing for ``build USD from source'': component control (strip
imaging on a farm), dependency pinning to match a DCC's runtime, and ABI
parity for in-house plug-ins. Know the tool's name ---
\texttt{build\_usd.py} --- and that prebuilt binaries (pip
\texttt{usd-core}, NVIDIA releases) trade those choices away for
convenience.
\end{ExamTip}

\subsection{Custom kinds}\label{sec:cu-kinds}\objref{3.3}
\index{kind}The kind vocabulary itself is extensible --- a
\texttt{plugInfo.json} resource registers new kinds into the hierarchy:

\begin{lstlisting}
{
    "Plugins": [{
        "Type": "resource",
        "Name": "studioKinds",
        "Info": {
            "Kinds": {
                "vehicle": { "baseKind": "component" }
            }
        }
    }]
}
\end{lstlisting}

\noindent\texttt{baseKind} slots the new kind into existing semantics ---
a \texttt{vehicle} \emph{is a} component, so model-hierarchy tools keep
working. Create custom kinds when tools must select your category
(``all vehicles''), and remember the registration must exist on
\emph{every} site that consumes the assets --- kinds are vocabulary, not
data.

Registration end to end is exactly two steps: put the
\texttt{plugInfo.json} in a directory, put that directory on
\texttt{PXR\_PLUGINPATH\_NAME}. Verification is three lines --- and this
book's test suite runs precisely this in a subprocess against the
\texttt{plugInfo.json} printed above:

\begin{lstlisting}[language=Python]
# PXR_PLUGINPATH_NAME=/studio/plugins/studio_kinds python
from pxr import Kind

print(Kind.Registry.HasKind("vehicle"))      # True
print(Kind.Registry.GetBaseKind("vehicle"))  # component
\end{lstlisting}

\noindent Codeless schemas ride the same mechanism --- their plugin
directory carries a \texttt{generatedSchema.usda} plus a
\texttt{plugInfo.json} with \texttt{Types} entries. Let
\texttt{usdGenSchema} generate both from your \texttt{schema.usda}:
the \texttt{Types} block encodes TfType registration details that are
easy to get subtly wrong by hand.

\subsection{Resolvers, file formats, and SceneIndex --- who does what}\label{sec:cu-ext}\objref{3.5, 3.7, 3.8}
Three extension points get conflated on exams; the boundaries are clean:

An \textbf{Ar resolver}\index{resolver} answers ``\emph{where} is asset
\texttt{X}?'' --- it maps logical identifiers
(\texttt{asset://props/chair?version=latest}) to concrete resolved paths,
honoring contexts (site, show, pin file). It never parses content.

A \textbf{file format plug-in}\index{file format plug-in} answers ``how do I
read \emph{this kind of file} as USD?'' --- it parses a foreign format into
scene description on demand (that is how \texttt{usda}/\texttt{usdc}/Alembic
themselves are wired in). Dynamic file formats can even \emph{generate}
prims from parameters.

A \textbf{SceneIndex plug-in}\index{SceneIndex} sits on the Hydra side:
it filters or procedurally generates \emph{render} prims after composition
--- geometry the stage never authored (the ``in-memory renderable
primitives'' objective lives here and in dynamic file formats).

\begin{Gotcha}
``Build a plug-in against a specific USD version'' is not ceremony: USD's
C++ ABI changes between releases, so a plug-in must be compiled against the
\emph{exact} USD build that loads it --- same version, same compiler
family, same flags. A plug-in built for a studio's USD 24.08 will not load
into the pip \texttt{usd-core} of a different version; the symptom is a
silent non-load, and the diagnosis is \texttt{TF\_DEBUG=PLUG\_REGISTRATION}.
\end{Gotcha}

\DrillsTitle
\subsection*{Drill 1: Typed vs API schema}
Give 2 examples of when you would choose each.
\AnswerLines{6}
\begin{DrillSolution}
Typed/IsA: a robot joint (\texttt{SimJoint}) or a studio-specific camera rig
--- entities with their own identity that prims \emph{are}. Applied API:
physics properties on existing meshes, wear/aging data on any prop ---
capabilities added to prims that already have a type. Heuristic: \emph{is}
vs \emph{has}.
\end{DrillSolution}

\subsection*{Drill 2: Resolver design}
Sketch an example: \texttt{asset://chair?variant=oak} $\rightarrow$ actual
file path.
\AnswerLines{7}
\begin{DrillSolution}
The resolver parses scheme + identifier + query; consults its context (show
config, pin file) for the published version; returns e.g.\
\texttt{/site/assets/chair/v012/chair\_oak.usdc}. Scene files keep the
logical URI; re-pinning the context retargets every scene at once. Cache
invalidation and farm/workstation context parity are the two operational
risks.
\end{DrillSolution}

\subsection*{Drill 3: Pick the mechanism}
Custom attribute, API schema, typed schema, or custom kind --- one per case:
(a) per-prim review notes; (b) ``all vehicles'' selection in tools; (c) a
new \texttt{ConveyorBelt} entity; (d) torque limits on existing joints
across shows.
\AnswerLines{6}
\begin{DrillSolution}
(a) namespaced custom attribute (or \texttt{customData} if untyped and
unsampled); (b) custom kind with \texttt{baseKind = component}; (c) typed
schema; (d) applied API schema --- shared, documented, validatable across
shows.
\end{DrillSolution}

\section*{Checkpoint Questions}
\addcontentsline{toc}{section}{Checkpoint Questions}

\noindent\textbf{CQ1.} A colleague's new schema ``does nothing'' on the farm
but works locally. First two checks?
\begin{DrillSolution}
Is the plugin discoverable there (\texttt{PXR\_PLUGINPATH\_NAME} includes
its \texttt{plugInfo.json}? \texttt{TF\_DEBUG=PLUG\_REGISTRATION}), and is
the farm's USD the same version/ABI the plug-in was built against? Codeless
schemas dodge the second problem entirely.
\end{DrillSolution}

\noindent\textbf{CQ2.} Why are codeless schemas the default recommendation
for pipeline data models?
\begin{DrillSolution}
They deploy as data: no compilation, no ABI coupling, no per-platform
builds --- just plugInfo + the generated schema layer on the plugin path.
You give up generated C++/Python accessor classes; generic
\texttt{UsdPrim}/\texttt{UsdAttribute} access and \texttt{HasAPI} still
work.
\end{DrillSolution}

\noindent\textbf{CQ3.} Match the objective to the extension point:
``generate renderable prims that were never authored on the stage.''
\begin{DrillSolution}
Hydra-side generation --- a SceneIndex plug-in (or, upstream of render, a
dynamic file format producing prims from parameters). An Ar resolver is the
wrong answer by definition: it locates content, it does not create it.
\end{DrillSolution}

\subsection*{Mini-checklist (tick when mastered)}
\begin{itemize}[leftmargin=*]
  \item[$\square$] I can explain what a resolver does in one sentence.
  \item[$\square$] I can justify a schema over ad-hoc attrs --- and codeless over compiled.
  \item[$\square$] I can place kinds/schemas/formats/resolvers/SceneIndex on the plugin map.
  \item[$\square$] I know why plug-ins pin exact USD versions.
\end{itemize}
